\scanpage{41}
b) \uline{Filtration}s\deleted{\unsure{ass}} \uline{associées aux cristaux de Dieudonné}. Revenons au cas d'un S général (sur lequel p soit loc. nilp.). On \unsure{définita|définira} alors, pour tout $G\in\operatorname{Ob}\mathrm{BT}(S)$, une filtration \deleted{en deux crans} de MoDule localement libre $\underline D^*(G)_S$ sur S par un sous-Module localement facteur direct (donc localement libre de t.f.) \inserted{$\uline{Fil}^{1}=$} $\underline\omega_G$, donnant lieu à une suite exacte
\[
(*)\qquad
\diagram{\begin{tikzcd}
0 \arrow[r] & \underline\omega_G \arrow[r] & \underline D^*(G) \arrow[r] & \underline t_{G^*} \arrow[r] & 0,
\end{tikzcd}}
\]
où en fait $\underline\omega_G$ est le module des différentielles invariantes sur le groupe formel $\overline G\subset G$ associé à G (qui sera construit plus bas), et $\underline t_G=\underline\omega_{G^*}^{\vee}$ est l'Algèbre de Lie associée au groupe de BT dual $G^*$ de G (définition à reporter aux généralités sur les goupes de BT). Cette \deleted{\unsure{filtr}} suite exacte est fonctorielle \inserted{en G} et compatible aux images inverses par un $S'\longrightarrow S$.

\margin{(*) \unsure{formule inutile si} $S_0$ \unsure{parfait}}
\margin{40}
